\section{\texorpdfstring{$[z,w]$ 的\personConley{}--\personZehnder{}指标}{[z,w] 的康利--曾德尔指标}}
\label{sec:conley-zehnder-index-capped-orbit}
\stfhSubIndexRangeStart{I065}{Kangli-Zengdeerzhibiao}{康利--曾德尔指标（Conley--Zehnder index）}{Zzwdezhibiao}{$[z,w]$ 的指标（of $[z,w]$）}

首先注意，辛群 $\operatorname{Sp}(S,\Omega)\cong\operatorname{Sp}(\mathbb R^{2n})$
有极大紧子群 $U(n)$，$\operatorname{Sp}(S,\Omega)$ 形变收缩到该子群
\footnote{译者注：原书在本段形变收缩句及下一页共轭作用句中，把同一辛向量空间的
形式从 $\Omega$ 改写成 $\omega$；译文按本节开头的定义统一为
$\operatorname{Sp}(S,\Omega)$。}。因此行列式映射
$\det:U(n)\to S^1$ 诱导自然同态
\[
 \mu:\pi_1(\operatorname{Sp}(S,\Omega),\mathrm{id})\longrightarrow\mathbb Z.
\]
接下来进一步展开第 2.2.2 节关于 $\operatorname{Sp}(S,\Omega)$ 结构的讨论。使用
(SZ92) 的记号，令
\[
 \operatorname{Sp}(S,\Omega)^*=
 \{A\in\operatorname{Sp}(S,\Omega)\mid\det(I-A)\neq0\}.
\]
注意，当 $p\in\operatorname{Fix}\phi_H^1$ 非退化时，
$d\phi_H^1(p)\in\operatorname{Sp}(T_pM,\omega_p)$。下面由 (CZ84) 证明的引理对后续
讨论十分重要；其证明留给原文或 (SZ92)。

\begin{lemma}
\label{lem:18-symplectic-group-star-components}
$\operatorname{Sp}(S,\Omega)^*$ 有两个连通分支
\[
 \operatorname{Sp}^{\pm}(S,\Omega)=
 \{A\in\operatorname{Sp}(S,\Omega)^*\mid
   \pm\det(I-A)>0\}.
\]
此外，$\operatorname{Sp}(S,\Omega)^*$ 中的每个环路在
$\operatorname{Sp}(S,\Omega)$ 中可缩。
\end{lemma}

因此，$\operatorname{Sp}(S,\Omega)^*$ 和
$\operatorname{Sp}(S,\Omega)^\pm$ 在 $\operatorname{Sp}(S,\Omega)$ 的共轭作用下不变。
考虑道路集
\begin{equation}
 \operatorname{SP}^*(1)=\{\alpha:[0,1]\to\operatorname{Sp}(S,\Omega)
 \mid\alpha(0)=\mathrm{id},\ \alpha(1)\in\operatorname{Sp}^*(S,\Omega)\}.
 \label{eq:18-sp-star-path-space}
\end{equation}
不失一般性，设 $\alpha(1)\in\operatorname{Sp}^+(2n;\mathbb R)$；
$\operatorname{Sp}^-(2n;\mathbb R)$ 的情形相同。由引理~\ref{lem:18-symplectic-group-star-components}，
可选取同伦意义下唯一的道路
$\alpha_+^{\mathrm{can}}:[1,2]\to\operatorname{Sp}(2n;\mathbb R)$，使
\begin{equation}
 \begin{gathered}
  \alpha_+^{\mathrm{can}}(1)=\alpha(1),\qquad
  \alpha_+^{\mathrm{can}}(2)=W^+,\\
  \alpha_+^{\mathrm{can}}(t)\in\operatorname{Sp}^+(2n;\mathbb R)
  \quad\text{对所有 }t\in[1,2],
 \end{gathered}
 \label{eq:18-canonical-extension-path}
\end{equation}
其中 $W^+=-\mathrm{id}$。
\footnote{译者注：原书把 $\alpha_+^{\mathrm{can}}$ 的定义域印作 $[0,1]$，但紧接着
指定其在 $t=1,2$ 的端值，并要求 $t\in[1,2]$；译文据此改为 $[1,2]$。}
对 $\alpha(1)\in\operatorname{Sp}^-(2n)$ 的道路，考虑矩阵
\[
 W^-=\operatorname{diag}(2,-1,\ldots,-1,1/2,\ldots,-1)
 \in\operatorname{Sp}^-(2n),
\]
并类似地定义 $\alpha_-^{\mathrm{can}}$。于是 (CZ84) 构造了下面的连续映射
$\rho:\operatorname{Sp}(V,\omega)\to S^1$；证明亦见 (SZ92) 的定理 3.1。

\begin{theorem}
\label{thm:18-salamon-zehnder-rho-map}
设 $(V,\omega)$ 是辛向量空间。则存在唯一的一族连续映射
\[
 \rho:\operatorname{Sp}(V,\omega)\longrightarrow S^1
\]
（每个 $(V,\omega)$ 一个），满足：
\begin{enumerate}[label=(\arabic*)]
\item \stfhPropertyLabel{（自然性）} 若 $T:(V_1,\omega_1)\to(V_2,\omega_2)$ 是辛同构，则对所有
$A\in\operatorname{Sp}(V_1,\omega_1)$，
$\rho(TAT^{-1})=\rho(A)$。
\item \stfhPropertyLabel{（乘积）} 若 $(V,\omega)=(V_1\times V_2,\omega_1\times\omega_2)$，
且 $A=A_1\oplus A_2\in\operatorname{Sp}(V,\omega)$，则
$\rho(A)=\rho(A_1)\rho(A_2)$。
\item \stfhPropertyLabel{（行列式）} 若
$A\in\operatorname{Sp}(2n,\mathbb R)\cap O(2n)=U_{\mathbb R}(n)$，则
$\rho(A)=\det(X+iY)$，其中
\[
 A=\begin{pmatrix}X&-Y\\Y&X\end{pmatrix},\qquad X+iY\in U(n).
\]
\item \stfhPropertyLabel{（归一化）} 若 $A$ 在单位圆上没有特征值，则 $\rho(A)=\pm1$。
\end{enumerate}
\end{theorem}

把 $\alpha$ 与 $\alpha_+^{\mathrm{can}}$ 连接，并考虑复合
\[
 \rho\circ(\alpha\#\alpha_+^{\mathrm{can}}):[0,2]\longrightarrow S^1,
\]
其端点值为 $\pm1$；因而将它平方后得到一个环路，其环绕数由映射
$\det^2:U(n)\to S^1$ 诱导。

\begin{definition}
\label{def:18-conley-zehnder-index-path}
设 $\alpha\in\operatorname{SP}^*(1)$。定义 $\alpha$ 的\personConley{}--\personZehnder{}指标为
\footnote{译者注：原书定义 18.4.3 印作
$\deg\det^2(\rho\circ(\alpha\#\alpha_+^{\mathrm{can}}))$；但 $\rho$ 的值已在
$S^1$ 中，不能再施加定义在 $U(n)$ 上的 $\det$。结合前句及 (SZ92) 的定义，译文把
它写成复合映射的平方之次数。}
\[
 \mu_{\mathrm{CZ}}(\alpha)
 =\deg\bigl((\rho\circ(\alpha\#\alpha_+^{\mathrm{can}}))^2\bigr).
\]
\end{definition}

定义~\ref{def:18-conley-zehnder-index-path} 的 $\mu_{\mathrm{CZ}}$ 与 (SZ92) 的定义
一致。另一方面，对给定\personLie{}群 $G$，记
\[
 \Omega(G)=\{g:S^1\to G\mid g(0)=\mathrm{id}\}
\]
为 $G$ 的基环路群。于是有明显的指标映射
$\mu:\Omega(\operatorname{Sp}(2n,\mathbb R))\to\mathbb Z$，它由形变收缩
$\operatorname{Sp}(2n,\mathbb R)\simeq U(n)$ 及复合映射的次数定义；具体地，对
$g\in\Omega(U(n))$，该整数由
\[
 g\longmapsto\frac1{2\pi}\int_{S^1}(\det\circ g)^*d\theta
\]
给出。
\footnote{译者注：原书在这里把相应同伦群印作
$\pi_1(\Omega(\operatorname{Sp}(2n,\mathbb R)),\mathrm{id})$ 与
$\pi_1(\Omega(U(n)),\mathrm{id})$，这会多取一层环路空间，且与紧接着对
$g:S^1\to\operatorname{Sp}(2n,\mathbb R)$ 的积分定义不相容；译文按后者陈述。}
\personConley{}--\personZehnder{}指标 $\mu_{\mathrm{CZ}}$ 可由下列性质刻画。

\begin{proposition}[{(FH93) 中的命题 5}]
\label{prop:18-conley-zehnder-characterization}
存在唯一映射
\[
 \mu_{\mathrm{CZ}}:\operatorname{SP}^*(1)\longrightarrow\mathbb Z
\]
满足以下准则：
\begin{enumerate}[label=(\arabic*)]
\item \stfhPropertyLabel{（归一化）}
\[
 \mu_{\mathrm{CZ}}\!\left(\tau\mapsto
 \{e^{-\tau}+ie^\tau\}\oplus\{e^{\pi i\tau}\}\oplus\cdots\oplus
 \{e^{\pi i\tau}\}\right)=n-1,
 \qquad
 \mu_{\mathrm{CZ}}(\tau\mapsto e^{\pi i\tau}\mathrm{Id})=n.
\]
\item \stfhPropertyLabel{（基环路作用）} 在基环路群
$\Omega(\operatorname{Sp}(2n,\mathbb R),\mathrm{id})$ 对
$\operatorname{SP}^*(1)$ 的作用下，对每个环路
$g:S^1\to\operatorname{Sp}(2n,\mathbb R)$，
\[
 \mu_{\mathrm{CZ}}(g\alpha)=2\mu(g)+\mu_{\mathrm{CZ}}(\alpha).
\]
\end{enumerate}
\end{proposition}

取 $[z,w]\in\operatorname{Crit}\mathcal A_H$，并选取\termTrivialization{}
\[
 \Psi:w^*TM\longrightarrow D^2\times\mathbb R^{2n}.
\]
回忆 $z(t)=\phi_H^t(p)$，其中 $p\in\operatorname{Fix}\phi_H^1$，于是可写
\[
 \Psi\circ d\phi_H^t\circ\Psi^{-1}(t,v)=(t,\alpha_\Psi(t)v).
\]
按定义，
\[
 \alpha_\Psi(0)=\mathrm{id},\qquad
 \alpha_\Psi(1)\in\operatorname{Sp}^*(2n;\mathbb R).
\]

\begin{definition}
\label{def:18-conley-zehnder-capped-orbit}
定义 $[z,w]$ 的\personConley{}--\personZehnder{}指标 $\mu_H([z,w])$ 为
\[
 \mu_H([z,w])=\mu_{\mathrm{CZ}}(\alpha_\Psi).
\]
\end{definition}

下面陈述“粘上一个球面”作用
$[z,w]\mapsto[z,w\#A]$ 下\personConley{}--\personZehnder{}指标与\termFirstChernNumber{}
$c_1(A)$ 之间的重要恒等式。其指标公式的证明沿用 (Oh06a) 的论述。

\begin{theorem}
\label{thm:18-conley-zehnder-capping-change}
设 $z:S^1=\mathbb R/\mathbb Z\to M$ 是
$\dot x=X_H(x)$ 的给定 1-周期解，$w,w'$ 是两个给定的填充圆盘。则
\begin{equation}
 \mu_H([z,w'])=\mu_H([z,w])-2c_1([w'\#\overline w]).
 \label{eq:18-index-two-cappings}
\end{equation}
特别地，
\begin{equation}
 \mu_H([z,w\#A])=\mu_H([z,w])-2c_1(A).
 \label{eq:18-index-sphere-gluing}
\end{equation}
\end{theorem}

\begin{remark}
\label{rem:18-conley-zehnder-sign-convention}
要强调的是，在我们的约定下，公式中第一\personChern{}数前的符号是“$-$”。它与 (HS95)
公式的符号不同，源于 $\mathbb C^n$ 上典范辛形式的约定不同：把
$\mathbb R^{2n}\cong T^*\mathbb R^n$，并把相应典范坐标记为
$(q_1,\ldots,q_n,p_1,\ldots,p_n)$ 时，我们约定的典范辛形式为
\begin{equation}
 \omega_0=\sum_jdq_j\wedge dp_j,
 \label{eq:18-canonical-symplectic-sign}
\end{equation}
而 (HS95)、(SZ92) 或 (Po01) 采用
\[
 \omega_0'=-\omega_0=\sum_jdp_j\wedge dq_j.
\]
\end{remark}

\begin{proof}
把指标公式的证明分成若干步。

\begin{enumerate}[label=\arabic*.]
\item \stfhPropertyLabel{（典范辛形式）} 在坐标 $z_j=q_j+ip_j$ 下，我们对
$T^*\mathbb R^n=\mathbb R^{2n}\cong\mathbb C^n$ 的典范辛形式采用
\eqref{eq:18-canonical-symplectic-sign} 的约定。

\item \stfhPropertyLabel{（典范复结构）} 设 $J_0$ 是 $\mathbb R^{2n}\cong\mathbb C^n$ 上由
复数 $i$ 的乘法给出的标准复结构，其中 $z_j=q_j+ip_j$。在我们对 $\mathbb C^n$ 上
典范辛形式 $\omega_0$ 的约定下，相伴的\personHermitian{}结构
\[
 \langle\cdot,\cdot\rangle:\mathbb C^n\times\mathbb C^n\longrightarrow\mathbb C
\]
\emph{对第一个变元复线性、对第二个变元反线性}。换言之，\personHermitian{}内积为
\begin{equation}
 \langle u,v\rangle=g(u,v)-i\omega_0(u,v),
 \label{eq:18-hermitian-our-convention}
\end{equation}
其中 $g$ 是 $\mathbb R^{2n}$ 上的标准欧氏内积。要注意，这个\personHermitian{}结构是
(HS95)、(SZ92)、(Po01) 所用结构的共轭；后者对应于
\begin{equation}
 \langle u,v\rangle=g(u,v)+i\omega_0(u,v).
 \label{eq:18-hermitian-other-convention}
\end{equation}
见 (SZ92) 引理 5.1 之前的注。等价地，后一个\personHermitian{}结构与殆\personKahler{}结构
\[
 (g,\omega_0',J_0')
\]
相伴，其中 $J_0'$ 是 $J_0$ 的共轭殆复结构。\emph{$\mathbb C^n$ 上复结构的这种改变，会
改变一般复向量丛 $E$ 的第一\personChern{}数的符号}；回忆复向量丛\personChern{}类的一般公式
（见 (MSt74) 引理 14.9）：
\[
 c_k(\overline E)=(-1)^kc_k(E).
\]

\item \stfhPropertyLabel{（$\operatorname{SP}^*(1)$ 上的\personConley{}--\personZehnder{}指标）}
注意，两种约定下 $\operatorname{Sp}(2n,\mathbb R)$ 的定义相同。

\item 给定两个映射 $w,w':D^2\to M$，满足
$w|_{\partial D^2}=w'|_{\partial D^2}$，如下定义粘合映射
$u=w\#\overline{w'}:S^2\to M$：
\[
 u(z)=\begin{cases}
  w(z),&z\in D^+,\\
  w'(1/\overline z),&z\in D^-.
 \end{cases}
\]
这里 $D^+$ 是保持原定向的 $D^2$，$D^-$ 带相反定向。先验地这只给出连续映射，
但通过在边界附近把映射“压平”，可在不改变同伦类的前提下将其形变为光滑映射
\footnote{译者注：原书说要把两个映射在边界附近“压平”，但紧接着的公式只写了
$w$。要使 $w\#\overline{w'}$ 沿粘合圆周光滑，两个圆盘都须作同样处理；译文补出
$w'$ 的条件。}。
换言之，可假设充分小的 $\epsilon>0$ 使得
\[
 \begin{gathered}
 w(z)=w(z/|z|),\qquad w'(z)=w'(z/|z|)\\
 \text{当 }|z|\geq1-\epsilon.
 \end{gathered}
\]
以下总假设填充圆盘在这个意义下是平坦的。作此调整后，$u$ 定义 $S^2$ 上的光滑映射。

\item \stfhPropertyLabel{（标记条件）} 对给定的 $[z,w]$、$[z,w']$，其中周期轨道
$z(t)=\phi_H^t(z(0))$，对平凡化
\[
 \Phi_w:w^*TM\to D^2\times(\mathbb R^{2n},\omega_0),\qquad
 \Phi_{w'}:(w')^*TM\to D^2\times(\mathbb R^{2n},\omega_0),
\]
附加总能满足的\emph{标记}条件
\begin{equation}
 \Phi_w(z(0))=\Phi_{w'}(z(0))
 \label{eq:18-trivialization-marking-condition}
\end{equation}
（视为从 $T_{z(0)}M$ 到 $\mathbb R^{2n}$ 的映射）。
\footnote{译者注：原书把 $\Phi_w$ 与 $\Phi_{w'}$ 的定义域都印作 $w^*TM$；后者按
其下标和后续使用应为 $(w')^*TM$，译文据此补正。}
\emph{在此附加条件下}，可以写成
\begin{equation}
 \alpha_{[z,w']}(t)=S_{w'w}(t)\cdot\alpha_{[z,w]}(t),
 \label{eq:18-capping-trivialization-loop}
\end{equation}
其中 $S_{w'w}:S^1=\mathbb R/\mathbb Z\to\operatorname{Sp}(2n,\mathbb R)$ 是由
\eqref{eq:18-capping-trivialization-loop} 定义的\emph{环路}。它确实是环路，因为
\begin{equation}
 \alpha_{[z,w']}(0)=\alpha_{[z,w]}(0)=\mathrm{id},
 \label{eq:18-capping-path-start}
\end{equation}
\begin{equation}
 \alpha_{[z,w']}(1)=\alpha_{[z,w]}(1).
 \label{eq:18-capping-path-end}
\end{equation}
事实上，由~\eqref{eq:18-capping-trivialization-loop} 与
\eqref{eq:18-trivialization-marking-condition}，有恒等式
\begin{align}
 S_{w'w}(t)
 &=\bigl(\Phi_{w'}(z(t))\circ d\phi_H^t(z(0))
          \circ\Phi_{w'}(z(0))^{-1}\bigr)\nonumber\\
 &\quad\circ\bigl(\Phi_w(z(t))\circ d\phi_H^t(z(0))
          \circ\Phi_w(z(0))^{-1}\bigr)^{-1}\nonumber\\
 &=\Phi_{w'}(z(t))\circ d\phi_H^t(z(0))\nonumber\\
 &\quad\circ\Phi_{w'}(z(0))^{-1}\circ\Phi_w(z(0))\nonumber\\
 &\quad\circ(d\phi_H^t)^{-1}(z(0))\circ(\Phi_w(z(t)))^{-1}.
 \label{eq:18-transition-loop-expanded}
\end{align}
\stfhIndex{I123}{Biaojitiaojian}{标记条件（marking condition）}
标记条件~\eqref{eq:18-trivialization-marking-condition} 使
\eqref{eq:18-transition-loop-expanded} 中间的项相消，因而
\begin{equation}
 S_{w'w}(t)=\Phi_{w'}(z(t))\circ\Phi_w(z(t))^{-1}.
 \label{eq:18-transition-loop}
\end{equation}
由命题~\ref{prop:18-conley-zehnder-characterization} 及
\eqref{eq:18-transition-loop}，得到
\begin{equation}
 \operatorname{ind}_{\mathrm{CZ}}(\overline\alpha_{[z,w']})
 =2\operatorname{wind}(\widehat S_{w'w})
  +\operatorname{ind}_{\mathrm{CZ}}(\overline\alpha_{[z,w]}).
 \label{eq:18-index-transition-winding}
\end{equation}
这里 $\widehat S_{w'w}:S^1\to U(n)$ 是在
$\operatorname{Sp}(2n,\mathbb R)$ 内与 $S_{w'w}$ 同伦的酉矩阵环路。这样的同伦总存在，
并且在同伦意义下唯一，因为 $U(n)$ 是 $\operatorname{Sp}(2n,\mathbb R)$ 的形变收缩。

\stfhIndex{I079}{Diyichenlei}{第一陈类（first Chern class）}
\item \stfhPropertyLabel{（$c_1$ 的归一化）} 最后回忆 $S^2$ 上辛向量丛 $E$ 的第一\personChern{}类
$c_1$ 的定义。我们把\personChern{}类归一化，使
$S^2\cong\mathbb{CP}^1$ 的切丛具有第一\personChern{}数 $2$；这也与文献中的标准约定一致。
\emph{要指出的是，这种归一化与我们约定中式~\eqref{eq:18-hermitian-our-convention}
给出的 $\mathbb C^n$ 上\personHermitian{}结构相容}（见 (Mil65) 第 167 页）。

分解 $S^2=D^+\cup D^-$，并考虑辛平凡化
\[
 \Phi_+:E|_{D^+}\to D^2\times(\mathbb R^{2n},\omega_0),\qquad
 \Phi_-:E|_{D^-}\to D^2\times(\mathbb R^{2n},\omega_0).
\]
在我们约定的 $\mathbb C^n$ 上\personHermitian{}结构下，它们同伦于酉平凡化；在另一约定
下，则同伦于共轭酉平凡化。记转移矩阵环路为
\[
 \phi_{+-}:S^1\longrightarrow\operatorname{Sp}(2n,\mathbb R),
\]
它由方程
\[
 \Phi_+|_{S^1}\circ(\Phi_-|_{S^1})^{-1}(t,\xi)
 =(t,\phi_{+-}(t)\xi)
\]
确定，其中 $(t,\xi)\in E|_{S^1}$，且
$S^1=\partial D^+=\partial D^-$。按定义，在我们的约定下
\begin{equation}
 c_1(E)=\mu(\widehat\phi_{+-}).
 \label{eq:18-chern-transition-our-convention}
\end{equation}
等价地，
\begin{equation}
 c_1(E)=-\mu(\overline{\widehat\phi}_{+-}).
 \label{eq:18-chern-transition-conjugate}
\end{equation}

现在把它应用到 $u^*(TM)$，其中 $u=w\#\overline{w'}$，$\Phi_w,\Phi_{w'}$ 是
第 5 步给出的平凡化。由~\eqref{eq:18-transition-loop}，$S_{w'w}$ 是
$\Phi_w$ 与 $\Phi_{w'}$ 之间的转移矩阵环路。因此按定义，第一\personChern{}数
$c_1(u^*TM)$ 就是酉矩阵环路
\[
 \widehat S_{w'w}:t\longmapsto\widehat S_{w'w}(t),
 \qquad S^1\longrightarrow U(n)
\]
在我们约定的 $\mathbb C^n$ 上\personHermitian{}结构中的环绕数
$\mu(\widehat S_{w'w})$。容易验证，把这个环绕数用于 $S^2$ 的切丛时，其值确为
$2$，因此与我们采用的\personChern{}类约定一致。
\footnote{译者注：原书此处称 $\Phi_w,\Phi_{w'}$ 在第 4 步给出；两者实际在第 5 步
“标记条件”中引入，译文据此更正步骤号。}

\item \stfhPropertyLabel{（证明收尾）} 综合以上各步，特别是第 2、5、6 步，公式
\eqref{eq:18-index-transition-winding} 和
\eqref{eq:18-chern-transition-conjugate} 就化为所要证明的指标公式。
\end{enumerate}
\end{proof}
